
\section{Introduction}\label{sec:intro}
A clearinghouse that accepts a basket of tokenised dollar claims as collateral has to answer a decision question that is standard in operational risk management: which of these instruments may be haircut together, and how much of the basket's joint tail is common? Answering it requires the joint dependence structure of the basket, and in practice that means a correlation matrix. The matrix is an input to a decision model, not an end in itself, but every downstream quantity - portfolio value at risk, expected shortfall, the margin multiplier applied to a correlated group - inherits whatever the matrix gets wrong.

The difficulty is that the matrix cannot be estimated from the data that the decision problem generates. On a representative day only three or four coins in a panel of dozens deviate from par at all; the rest supply a flat line of near-zero deviations. A sample correlation computed over such a window is dominated by the flatness, and the matrices that clearinghouses, index providers and regulators publish are accordingly smoothed across the regions where the deviations are missing. The smoothing is defended as regularisation under sparse data, and within a single regime that defence is sound.

It fails across regimes. A stablecoin's peg survives or fails according to structural conditions that change discretely: whether collateral backs it at par, whether its price feed is still reporting, whether redemption is open or gated, whether the custodian holding its reserves is solvent, whether a bridge has been compromised, whether a regulator has intervened, whether the basis on the funding leg has compressed, and whether a buffer still absorbs the spread. Each is a channel through which a basket can be stressed, each has a threshold, and each defines a region of the parameter space within which the joint co-movement structure of the basket does not move at all. Across a threshold it steps. A smoother that is continuous in the parameters will therefore place at the boundary a value the true dependence never takes, and no quantity of additional within-regime data improves the estimate there.

This paper treats the resulting problem as one of testing for block structure in pairwise dependence when the block is prescribed ex ante rather than estimated. Fix an episode and the failure channel that triggered it; the channel identifies a set of coins, the \emph{wall-affected block}, independently of any observed co-movement. Classify every pair of coins by how many of its members lie in the block, and ask whether the change in correlation across the episode is largest on the block layer. The design is a test rather than a fit: the block is prescribed, the statistic is fixed before the data are seen, and the question is whether the realised pattern matches.

Three features of the procedure are worth stating because they are where the operational content lies. First, a crisis raises co-movement across the whole panel and not only inside the block, so the comparison that carries the inference is the block-versus-rest contrast, which differences out a component common to all pairs. Second, and less obviously, the unit of resampling must be the coin rather than the pair: pairs that share a member are not independent, so permuting pairs treats a dependent sample as if it were independent and overstates precision. Permuting instead over the coin sets that could have been labelled a block asks how unusual the realised block is among the blocks that the action could have taken, which is the question a decision maker actually faces. Third, the test does not require the amplitude of the crossing, which is left free, because the normalising map that takes covariances to correlations destroys scale while preserving sign and support.

We make three assumptions explicit rather than deriving them, so that what the procedure guarantees is separable from what it claims. Within a failure regime the co-movement covariance is constant up to a positive scale; across a regime boundary it jumps by a signed sum of rank-one terms; and the terms are supported on the wall-affected block. Everything downstream is proved: that correlation rigidity follows from covariance rigidity with the scale cancelled exactly, that the induced correlation jump has zero diagonal, that its image on the correlation matrix is a signed footprint in three tiers, and that no locally continuous estimator can track it. We then test the resulting prediction on the nine documented de-peg episodes of 2020--2025, using public price data only.

On the largest custodial failure in the record the block layer stands above both others, the block-versus-rest contrast is 0.69, and the coin-level permutation places the true block above 98.9\% of the 816 coin sets of the same size, giving $p=0.011$. The ordering holds across all nine specifications of a window grid. The pair formed by the two decentralised coins that contain neither the failed token nor any coin with direct custodial exposure shows the largest jump of any pair in the panel, which rules out the reading that the effect is carried by the failing instrument alone. Of the two further episodes with a prescribed multi-coin block, one is consistent in sign but imprecise and the other shows no block signal; we report them alongside rather than pooling them.

The remainder of the paper is organised as follows. Section~\ref{sec:moduli} defines the state space and states which properties are assumed and which derived. Section~\ref{sec:rigidity} develops the covariance and its normalised image and proves rigidity. Section~\ref{sec:wall} states the crossing law and derives its consequences for correlations. Section~\ref{sec:empirical} builds and runs the test. Section~\ref{sec:conclusion} concludes.
